%!TEX root = ../main.tex

\begin{abstract}

Autonomous navigation pipelines are usually built around a single platform:
replacing a sensor, or moving from a rover to a multirotor, means rewriting
perception and control. This thesis addresses the last link of a
hardware-independent ROS~2 autonomy pipeline --- deciding the motion --- with a
learnt policy that maps a sensor-agnostic voxel representation of the
surroundings, together with a declaration of the vehicle's kinematic limits,
directly to a velocity command. Because the local map forgets, data arrive at
irregular intervals and the right manoeuvre depends on what has just happened,
the policy is built on Liquid Neural Networks, continuous-time recurrent networks
whose state evolves with the interval actually elapsed.

Three cells --- LTC, CfC and LRCU --- each with a dense and a sparse Neural
Circuit Policy wiring, are implemented behind a sparse convolutional encoder and
compared, with everything else held equal, against a stateless control. Training
has two stages: imitation of a corpus of recorded flights, then fine-tuning with
a recurrent PPO on a fleet of parallel simulations of multirotors and rovers,
both instrumented with metrics designed to separate genuine learning from its
imitations.

A single policy commands both platforms by reading their limits from the
observation. Imitation generalises to unseen episodes with a test error about a
third of that of a velocity-echo baseline, and fine-tuning raises the share of
episodes that reach the goal from below 10\% to between 36\% and 63\%. No cell,
however, emerges as clearly superior, and the stateless control is among the best
in every stage, so the benefit of a time-consistent memory remains unproven; CfC
is indicated as the practical candidate for its fast convergence and its arrival
rate. The work also shows that the imitation loss does not predict closed-loop
performance, whereas the policy's dependence on the map is the imitation-stage
quantity most closely related to it.

\end{abstract}
